% P03.1 global-section-palette descent, live integration version.

\begin{definition}
\label{def:p031-global-section-palette}
Let $\pi:E\to X$ be a finite surjective morphism of $K$-analytic spaces.
A \emph{global section palette} for $\pi$ is a nonempty finite family of
analytic sections
\[
                 \sigma_i:X\longrightarrow E
                    \qquad(1\leq i\leq N)
\]
such that, for every $x\in X$,
\begin{equation}
\label{eq:p031-global-section-palette}
             \pi^{-1}(x)
                =\{\sigma_1(x),\ldots,\sigma_N(x)\}.
\end{equation}
Repetitions among the displayed points are allowed.
\end{definition}

\begin{theorem}
\label{thm:p031-section-palette-envelope-descent}
Let $K$ be complete and nontrivially valued, and let $\pi:E\to X$ be a
finite surjective morphism of good strict $K$-analytic spaces admitting a
global section palette.  Fix an obstacle
$f:X\to\mathbf R\cup\{-\infty\}$ and a semipositive real-PL cover
competitor $u\leq f\circ\pi$.  For every $x\in X$, there is a semipositive
real-PL base competitor $b_{x,u}\leq f$ such that
\begin{equation}
\label{eq:p031-section-palette-exact-repair}
 (b_{x,u}\circ\pi)(e)
   =\max_{e'\in\pi^{-1}(x)}u(e')
       \qquad(e\in\pi^{-1}(x)).
\end{equation}
The competitor $b_{x,u}$ may depend on $x$ and $u$.

Consequently, including when both competitor families are empty,
\begin{equation}
\label{eq:p031-section-palette-envelope-pullback}
                 s_{\mathrm{PL},E}
                    =s_{\mathrm{PL},X}\circ\pi.
\end{equation}
The same identity holds for the raw P03 envelope over the union of locally
Fubini--Study and semipositive real-PL minorants.
\end{theorem}

\begin{corollary}
\label{cor:p031-sectioned-orbit-defect-vanishing}
Keep the finite analytic orbit presentation of
\Cref{def:p031-finite-analytic-orbit-presentation}.  If $\pi$ admits one
global analytic section $\sigma:X\to E$, then
$\{g\circ\sigma:g\in G\}$ is a global section palette.  Hence
\eqref{eq:p031-section-palette-envelope-pullback} and the full P03 envelope
identity hold for every obstacle.  If $\mathcal P_X(f)$ is nonempty, then
the symmetry-breaking defect from
\Cref{thm:p031-centered-symmetry-defect} vanishes identically.
\end{corollary}

\begin{corollary}
\label{cor:p031-fold-envelope-without-maximum-closure}
In the fold-map example from
\Cref{prop:p031-real-pl-fiber-maximum-failure}, write
$q:X\amalg X\to X$ for the fold map, let $h$ be the semipositive real-PL
function on $X$, and set $f:=\max\{h,0\}$.  Let $u$ equal $h$ on the first
component and $0$ on the second.  Then
\begin{equation}
\label{eq:p031-fold-envelope-values}
 s_{\mathrm{PL},X}=f,
 \qquad
 s_{\mathrm{PL},X\amalg X}=f\circ q,
 \qquad
 \Sigma_{q,f}\equiv0,
\end{equation}
although
\[
                         q_{\max}u=f
\]
is not real PL.  Thus failure of real-PL fiber-maximum closure does not by
itself produce a nonzero envelope defect.
\end{corollary}
